\chapter{Introduction}
\label{ch:intro}

\section{Motivation}

Vision transformers pretrained on large image collections are usually adapted to a new task by fine-tuning.
\emph{Full fine-tuning} updates every weight of the model. For a model of the size studied here (about 86 million
weights) this has three costs: training needs several gibibytes of GPU memory, because the intermediate results of
every layer must be kept for the backward pass; every task produces a complete copy of the model, about 330~MiB; and
a gradient is computed for every weight.

\emph{Low-rank adaptation} (LoRA)~\cite{hu2022lora} is the most widely used alternative. It keeps the pretrained
weights frozen and trains a small low-rank correction next to selected weight matrices, together with a new
classifier. The trained correction can afterwards be added into the weight matrix (\emph{merged}), which gives a
model of exactly the original form. LoRA was proposed for language models, and most of what is known about its cost
comes from there. For on-device use, where a model is adapted on a device with little memory, the relevant questions
are practical: how much memory and time does LoRA really save, at what cost in accuracy, and does it change anything
once the model is deployed?

This project answers these questions by measurement for two vision transformers of different design. ViT-B/16
\cite{dosovitskiy2021vit} is a plain transformer in which all twelve blocks have the same size. Swin-B
\cite{liu2021swin} is hierarchical: its 24 blocks work in four stages at decreasing resolution. Both are fine-tuned
on CIFAR-100~\cite{krizhevsky2009cifar}. Every configuration of the main comparison is run on both models with the
same settings and the same five seeds, so that a statement about LoRA can be checked on two architectures.

\section{Questions and findings}

\begin{enumerate}
  \item \textbf{Accuracy: is LoRA as accurate as full fine-tuning?} With the learning rate of each method chosen on
        validation loss, LoRA at rank 8 is more accurate than full fine-tuning on Swin-B, by 0.4--0.5 percentage
        points (pp) for each of the three targets. On ViT-B there is no evidence that it is better: two targets are level with
        full fine-tuning, and attention LoRA is 0.2~pp below. The comparison depends on tuning full fine-tuning: at a learning rate three
        times smaller, full fine-tuning of ViT-B is 0.6~pp less accurate. Rank 16 gives no gain over rank 8 on
        ViT-B (Chapter~\ref{ch:lora_ft}).
  \item \textbf{Training cost: what does LoRA save?} The saving in training memory depends on which layers are
        adapted, in the same way on both models. At batch size 64, adapters on the attention projections save
        about a third of the peak memory, adapters on the MLP about a seventh, and adapters on both about 6\%. The
        saving is largest for small batches: about 70\% at batch size 1. A training epoch is 12--26\% shorter
        (Chapter~\ref{ch:lora_ft}).
  \item \textbf{Placement: which blocks should receive adapters?} Training memory is set mainly by the first block
        that has an adapter, because the backward pass has to reach it. Adapting only the second half of the blocks
        therefore needs the least memory, but it loses 1.2--1.6~pp and falls below full fine-tuning on both models.
        Adapting every second block and skipping the first one (the odd blocks) saves 11--14\% of the memory; on Swin-B
        its accuracy is level with adapters on all blocks, on ViT-B a small loss cannot be excluded (Chapter~\ref{ch:placement}).
  \item \textbf{Deployment: what changes at prediction time?} For a single task, nothing, provided the adapters are
        merged: a merged model gives the same predictions and has the memory and speed of an ordinary model.
        Adapters kept separate make single-image predictions 8--55\% slower on the GPU. The small stored file
        (2--6~MiB per task with adapters on all blocks) is an advantage only when several tasks share one base model
        (Chapter~\ref{ch:deploy}).
\end{enumerate}
A fifth question was added towards the end of the project: given a limit on training memory, can the placement be
chosen in advance? Peak training memory turned out to be predictable from the shape of the model (a formula fitted
on ViT-B predicted ViT-L to within 1.5\% before it was measured), but a simple rule for the accuracy of a
placement, written down before the runs, failed (Chapter~\ref{ch:budget}).

\section{Scope and organisation}

All results are for one dataset (CIFAR-100, with the images upsampled to the input size of the models), in 32-bit
floating point, on one partition of a shared server GPU. No edge device was used, and the power drawn by a
single job could not be measured. Chapter~\ref{ch:background} gives the background and related work, and
Chapter~\ref{ch:setup} the experimental setup and the definitions of all metrics.
Chapters~\ref{ch:lora_ft} to~\ref{ch:budget} present the results, and Chapter~\ref{ch:conclusion} concludes with
limitations and future work. The appendices hold the detailed tables.
